\documentclass[../main]{subfiles}
\begin{document}

\section{Programa}

\begin{enumerate}
  \item HERRAMIENTAS y CONOCIMIENTOS PREVIOS
    \begin{itemize}
        Que el alumno Sepa seleccionar e implementar herramientas de
        planificación como:
      \item Diagramas de Gantt.
      \item Pizarras o tarjetas kamban
      \item Diagrama iterativo de toma de decisiones
      \item Diagrama espina de pescado
      \item Diagrama de flujo
      \item Diagrama de árbol (toma de decisiones
        \item Estudio de tiempos y movimientos
        \item Nuevos Materiales
        \item Principios para generar una FEM
        \item Nanotecnología
        \item Diferentes Efectos Físicos y Químicos para
          generar energías.
      \end{itemize}

    \item SELECCIÓN DEL PROYECTO
      \begin{itemize}
        \item Después de explorar herramientas y conceptos
          adicionales el alumno tendrá que realizar la
          elección de tres posibles proyectos del ámbito
          electromecánico que involucre la fabricación de
          algún producto, la generación de energías y
          otro con fines solidarios.
        \item La elección del proyecto (viabilidad -
          Beneficio)se ará en acuerdo al mejor criterio
          Alumno -profesor
      \end{itemize}

    \item ETAPAS DEL PROYECTO
      \begin{itemize}
        \item Imagen del producto (Logotipo, packaging,
          colores)
        \item Constitución de la empresa S.A., S.R.L.,
          cooperativa etc. Rubro.
        \item Descripción del producto o servicio, como se va
          a realizar o cuales son las partes del proceso productivo.
        \item Objetivos Generales y Específicos
        \item Análisis de la concepción del proyecto, cómo
          surgió la idea, que necesidades encontraron o
          que satisfacción quieren lograr.
        \item Análisis de Mercado, oferta y demanda,
          competencia, productos sustitutos y
          complementarios.
        \item Marco Jurídico, análisis de costos
        \item Diagramas necesarios para la realización del
          proyecto, de gant, flujo, planos etc.
        \item Plan de seguridad e Higiene
        \item Manual de usuario, mantenimiento
        \item Marketing (página web, video promocional etc.
          \item Realización de Maquetas a pequeña escala
        \end{itemize}
    \end{enumerate}
    %\tableofcontents
    \end{document}

